\chapter{The Media System}
\label{ch:media-system}

\section{A correction worth stating up front}

Volume~I's opening chapter quoted CNA's own README verbatim: the \cnans{Media} namespace is
``25/25 present, but 14 are shells.'' That figure is real, but it is also, at the time of
writing, \textbf{stale} --- confirmed by directly reading the namespace's own current
tracking record. A dedicated, from-scratch effort ran sixteen phases and 231 tasks against
this namespace, closing 224 of them, backed by fifteen independent adversarial external
review rounds. All 24 \cnans{Media} types are now real implementations backed by genuine file
parsing --- Vorbis and ID3v2 tag parsing, \texttt{.m3u}/\texttt{.m3u8} playlist parsing,
FFmpeg-backed video, SDL3\_mixer-backed audio --- not shells. 5{,}471 tests exist for this
namespace, with 5{,}467 passing and only four hardware-only skips (accelerometer/gyroscope).
This is exactly the kind of documentation drift Volume~I's first chapter warned readers to
watch for: a specific, dated, checkable claim that this book preserves as evidence of the
project's own earlier state, corrected here against the namespace's own current record rather
than silently repeated as still current.

\section{Song, and the one gap even FNA has}

\cnaclass{Song} is a real, file-backed asset, constructed either from an explicit path or
discovered through a library scan (\S\ref{sec:medialibrary}). \texttt{Album}/\texttt{Artist}/%
\texttt{Genre} back-references are non-owning and remain null unless the song was actually
discovered by a scan rather than constructed directly. \texttt{IsRated}/\texttt{Rating}
(on XNA's 0--10 scale) are populated from real ID3v2 \texttt{POPM} or Vorbis \texttt{RATING}
tags, with \texttt{IsRated} deliberately reporting false for an explicit rating of zero,
since both tag formats reserve that value to mean ``unrated'' rather than ``rated zero.''
\texttt{IsProtected} is always false --- correctly, since a DRM-wrapped file is not something
a local library scan could index in the first place.

One finding here is worth its own paragraph, because of what it reveals about verifying
against the wrong reference: FNA's own \texttt{Song} class omits \texttt{Album},
\texttt{Artist}, \texttt{Genre}, and \texttt{ToString()} relative to the real XNA~4.0
reference assemblies --- a gap invisible across eight separate audit rounds, because every
one of them diffed against FNA's source rather than against Microsoft's own original
reference documentation directly. It was only found by finally diffing against the real
XNA~4.0 API surface itself (the same \texttt{xna4-spec} resource covered in
Chapter~\ref{ch:xna4-spec-auditing}). This is the sharpest illustration in this entire book of
why Volume~I's first chapter insisted on distinguishing ``matches FNA'' from ``matches real
XNA'' as two different claims: here, they were not the same claim at all, and treating FNA as
an infallible stand-in for the real spec cost eight rounds of review before the gap surfaced.

\section{The library hierarchy}
\label{sec:medialibrary}

\cnaclass{SongCollection}, \cnaclass{Album}/\texttt{AlbumCollection},
\cnaclass{Artist}/\texttt{ArtistCollection}, and \cnaclass{Genre}/\texttt{GenreCollection}
form a real, populated hierarchy, built by \cnaclass{MediaLibrary} through a genuine
filesystem scan of the host's music and pictures roots. \texttt{Album::GetAlbumArt()}/%
\texttt{GetThumbnail()} return a real \texttt{System::IO::Stream}, resolved either by
filename convention (\texttt{cover.jpg}, \texttt{folder.jpg}) or extracted on demand from an
embedded ID3v2 \texttt{APIC} or FLAC \texttt{PICTURE} frame --- extraction directly from an
embedded frame at scan time, rather than on demand, is a documented, deliberate future
increment, not a bug. \cnaclass{Picture}/\texttt{PictureAlbum} and their collections form a
parallel, real hierarchical tree for image libraries; \texttt{MediaLibrary::SavePicture()}
lazily creates a real ``Saved Pictures'' node the first time it is needed.

\section{Playback: MediaPlayer, MediaQueue, Playlist}

\cnaclass{MediaPlayer} is the static playback controller familiar from real XNA:
\texttt{Play(Song*)}/\texttt{Play(SongCollection)}, \texttt{Pause}/\texttt{Resume}/%
\texttt{Stop}, a real \cnaclass{MediaQueue} with shuffle and repeat, and
\texttt{GetVisualizationData()}. Playback position tracking is real, timer-based state, with
an explicit fallback song-end detector for cases where no native ``track finished'' signal is
available from the underlying decoder. \cnaclass{Playlist}/\texttt{PlaylistCollection}
support genuine, read-only \texttt{.m3u}/\texttt{.m3u8} parsing --- there is no write API, but
that matches real XNA's own \cnaclass{Playlist} being read-only too, not a CNA-specific
limitation.

\section{Video and VideoPlayer}

\cnaclass{Video}/\texttt{VideoPlayer} are real, backed by genuine FFmpeg decoding, CNA's own
graphics-backend frame rendering, and SDL3 audio streaming, on Linux and macOS.
\texttt{VideoPlayer::Play()} deliberately throws \texttt{InvalidOperationException} if a
video's own declared width, height, or frame-rate metadata does not match what the decoded
file itself actually reports --- a real, checked consistency guard rather than a silent
mismatch. \cnaclass{VideoSoundtrackType} (\texttt{Music}/\texttt{Dialog}/%
\texttt{MusicAndDialog}) is confirmed metadata-only in both FNA and CNA alike; neither
implements any actual audio-ducking or muting behavior tied to this value.

\section{The one genuinely open gap: Windows, Android, and Emscripten video}

The one area of this namespace still genuinely open, rather than merely undocumented, is
real FFmpeg-backed video on Windows, Android, and Emscripten --- excluded from those builds
entirely via a \texttt{CNA\_FFMPEG\_AVAILABLE} CMake gate. Using \texttt{Video}/%
\texttt{VideoPlayer} there today is a \emph{link error}, not a runtime stub throwing a
graceful exception --- a direct consequence of the project owner explicitly rejecting a
\texttt{NotSupportedException}-stub approach in favor of holding out for a real FFmpeg
integration on those platforms instead. This is not closeable from a Linux-only development
environment, and is recorded here as exactly that kind of gap: not silently missing, not
falsely claimed complete, but named and left for real cross-platform FFmpeg packaging work.

\section{Deliberate, documented non-gaps}

A short list of choices that look like gaps but are not: \texttt{.m4a}/\texttt{.aac}/WMA
files are deliberately not indexed by \cnaclass{MediaLibrary} at all, because SDL3\_mixer
ships no AAC decoder --- indexing a song the player could never actually play would be worse
than not listing it. There is no \texttt{MediaLibrary::Refresh()}, because a library scan is
a deliberate one-time synchronous snapshot, not a live-updating view. And
\texttt{MediaPlayer::Update()}'s fallback branch for a build with sound disabled is,
confirmed directly, dead code in every build configuration this project currently produces
--- untestable as a consequence of that, and not a Media-specific concern.
